% Abstract for "Unown - Pokemon Legendary Status Predictor"
\chapter*{Abstract}

This project addresses automated identification of legendary Pokémon using structured game-derived attributes and domain-informed feature engineering. Motivation stems from a highly imbalanced classification problem where a small subset of Pokémon exhibit ``legendary'' status; automated detection can aid dataset curation, analytics, and interactive tools. We use a curated dataset compiled into \texttt{pokemon\_data.csv} and apply a bespoke \texttt{FeatureBuilder} to extract priors (type- and ability-based legendary rates), body-metrics, z-scored statistics, and composite indices that encode offensive/defensive profiles.

A range of supervised classifiers were evaluated; the production solution uses a Random Forest trained on engineered features, with hyperparameter tuning and an optimized decision threshold to balance precision and recall for the rare class. The Random Forest achieved strong discrimination (AUC $\approx$ 0.97) and, after threshold optimization, an F1 $\approx$ 0.78 on the test fold, demonstrating robust recall/precision trade-offs for legendary detection. Model artifacts (trained model, feature list, and the feature builder) are serialized to disk (\texttt{pokemon\_legendary\_model.joblib}, \texttt{feature\_columns.joblib}, \texttt{feature\_builder.joblib}) and a Streamlit app provides an interactive demo.

Overall, the approach shows that combining simple domain priors (type/ability frequencies) with carefully constructed statistical features yields a compact, explainable model that reliably highlights legendary Pokémon in an imbalanced setting.
\chapter*{Abstract}

This project addresses automated identification of legendary Pokémon using structured game-derived attributes and domain-informed feature engineering. Motivation stems from a highly imbalanced classification problem where a small subset of Pokémon exhibit ``legendary'' status; automated detection can aid dataset curation, analytics, and interactive tools. We use a curated dataset compiled into `pokemon_data.csv` and apply a bespoke `FeatureBuilder` to extract priors (type- and ability-based legendary rates), body-metrics, z-scored statistics, and composite indices that encode offensive/defensive profiles.

A range of supervised classifiers were evaluated; the production solution uses a Random Forest trained on engineered features, with hyperparameter tuning and an optimized decision threshold to balance precision and recall for the rare class. The Random Forest achieved strong discrimination (AUC ≈ 0.97) and, after threshold optimization, an F1 ≈ 0.78 on the test fold, demonstrating robust recall/precision trade-offs for legendary detection. Model artifacts (trained model, feature list, and the feature builder) are serialized to disk (`pokemon_legendary_model.joblib`, `feature_columns.joblib`, `feature_builder.joblib`) and a Streamlit app provides an interactive demo.

Overall, the approach shows that combining simple domain priors (type/ability frequencies) with carefully constructed statistical features yields a compact, explainable model that reliably highlights legendary Pokémon in an imbalanced setting.
